% Generated by uscribe -- do not edit by hand
\documentclass[letterpaper,12pt]{article}
\usepackage[T1]{fontenc}
\usepackage[utf8]{inputenc}
\usepackage{graphicx}
\usepackage{hyperref}
\usepackage{xcolor}
\usepackage{fancyvrb}
\usepackage{booktabs}
\usepackage{longtable}
\IfFileExists{multirow.sty}{\usepackage{multirow}}{}
\hypersetup{colorlinks=true,linkcolor=blue,urlcolor=blue}
\begin{document}
\begin{titlepage}
\null\vfil\vskip 40pt
\begin{center}
{\LARGE\bfseries The Unicon Preprocessor \par}
\vskip 3em
{\large\bfseries Jafar Al-Gharaibeh \par}
\vskip 1.0em
{\bfseries Unicon Technical Report: 24\par}
\vskip 1.0em
{\large\bfseries April 2026\par}
\vskip 5.5em
{\bfseries Abstract}\par
\begin{quote}
The Unicon compiler runs a line-oriented source preprocessor before lexical analysis. The implementation, uni/unicon/preproce.icn, comes from Bob Alexander's Jcon preprocessor. It has been extended with triple-quoted multiline strings, function-like \$define macros, and built-in assert / assert\_not. This report describes those extensions, the architecture and directives they sit on, how \#line and includes interact with them, and how to add further facilities. Keywords: Unicon, preprocessor, triple-quoted strings, function-like macros, assertions, technical report.
\end{quote}
\vskip 0.8in
{\large\bfseries
Unicon Project\\
\url{http://unicon.org}\\
\ \
University of Idaho\\
Department of Computer Science\\
Moscow, ID, 83844, USA
}
\end{center}
\vfil\null
\end{titlepage}
\setcounter{page}{1}
\section{1. Introduction}
\label{utr24-sec-intro}

This report describes the Unicon source preprocessor: how a file plus includes becomes the token stream the compiler scans, what the \texttt{\$} directives do, and the three extensions added in \texttt{uni/unicon/preproce.icn} -- triple-quoted strings, function-like macros, and built-in \texttt{assert} / \texttt{assert\_not}. It is formatted as a Unicon Technical Report [1]. The language-facing command list lives in \emph{Programming with Unicon} (language reference, "Preprocessor") and UTR \#8 [2]; that text still states that \texttt{\$define} does not take arguments. Automated tests live under \texttt{tests/unicon/} (\texttt{triple\_strings.icn}, \texttt{macros*.icn}, \texttt{assert\_*.icn}). \texttt{unicon~-E} writes preprocessed source.


The preprocessor is not optional. Every compilation runs it.

\section{2. Motivation and prior art}
\label{utr24-sec-motivation}

The implementation is Bob Alexander's Jcon preprocessor, brought into Unicon with small edits [3]. It is line-oriented: it recognizes directives and identifiers, respects string and comment boundaries, and performs textual substitution. It is not a parser.


The Jcon baseline provided object-like \texttt{\$define}, \texttt{\$include}, \texttt{\$ifdef} / \texttt{\$ifndef} / \texttt{\$else} / \texttt{\$endif}, \texttt{\$error}, \texttt{\$line} / \texttt{\#line}, and predefined symbols from \texttt{\&features}. The implementation-book section "The Unicon Preprocessor" [4] documents that baseline (generator interface, include and \texttt{\$ifdef} stacks, quoted-string skipping, underscore continuation). It still says macros have no parameters. The language reference says the same [2].


Three gaps in that baseline are filled here:


\textbf{Multiline string literals.} Unicon string literals are one physical line unless continued with a trailing \texttt{\_}. A triple-quoted form lets a literal span lines, including text that looks like comments or \texttt{\$} directives, without changing the lexer grammar.


\textbf{Parameterized macros.} Object-like \texttt{\$define} cannot abstract repeated call-shaped text. Function-like \texttt{\$define} adds a parameter list and a replacement body, with C-style \texttt{\textbackslash{}} continuation on the definition and call arguments that may span lines.


\textbf{Assertions.} Debug and test builds need a check that disappears entirely in release. Built-in \texttt{assert} / \texttt{assert\_not} expand to runtime tests under \texttt{\_\_debug\_\_} / \texttt{\_\_test\_\_}, and to nothing otherwise.

\section{3. Architecture}
\label{utr24-sec-arch}

Canonical source is \texttt{uni/unicon/preproce.icn}. \texttt{uni/parser/preproce.icn} is a shorter, separate copy used by the standalone parser toolkit; it does not include the extensions in this report.

\subsection{3.1 Public procedures}
\label{utr24-sec-public}
\begin{longtable}{ll}
\hline
Procedure & Role \\
\hline
\endfirsthead
\hline
Procedure & Role \\
\hline
\endhead
\hline
\endfoot
\hline
\endlastfoot
\texttt{preprocessor(fname,~predefined\_syms)} & Main entry. Initializes state, loops over input, emits \\
\texttt{preproc(dummy,~args)} & Wrapper: \texttt{suspend~preprocessor(args[1],~predefs())}. \\
\texttt{predefs()} & Table of predefined symbols (\texttt{\_UNIX}, \texttt{\_\_DATE\_\_}, \\
\end{longtable}


\texttt{preproc\_new()} initializes; \texttt{preproc\_done()} clears large tables so they are not retained across compilations. \texttt{preproc\_install\_builtin\_macros()} runs from \texttt{preproc\_new()}. If \texttt{\_\_test\_\_} is enabled there, \texttt{\_\_debug\_\_} is defaulted.

\subsection{3.2 State}
\label{utr24-sec-state}

\texttt{record~preproc\_fmacro(params,~body)} holds a function-like macro: parameter names in order and replacement text. Object-like \texttt{\$define} symbols live only in \texttt{preproc\_sym\_table}.

\begin{longtable}{ll}
\hline
Global & Purpose \\
\hline
\endfirsthead
\hline
Global & Purpose \\
\hline
\endhead
\hline
\endfoot
\hline
\endlastfoot
\texttt{preproc\_sym\_table} & Object-like macros: symbol to replacement string. \\
\texttt{preproc\_fun\_table} & Function-like macros: symbol to \texttt{preproc\_fmacro}. \\
\texttt{preproc\_builtin\_fun\_table} & Names expanded specially (\texttt{assert}, \texttt{assert\_not}). \\
\texttt{preproc\_if\_stack}, \texttt{preproc\_if\_state} & Conditional compilation. \\
\texttt{preproc\_file\_stack}, \texttt{preproc\_file}, \texttt{preproc\_filename}, \texttt{preproc\_line} & Include stack and current input position. \\
\texttt{preproc\_include\_set}, \texttt{preproc\_include\_name} & Included paths; circular includes are errors. \\
\texttt{preproc\_print\_line}, \texttt{preproc\_print\_filename} & Last emitted position, for \texttt{\#line}. \\
\texttt{preproc\_ml\_anchor} & First physical line of a multiline macro call. \\
\texttt{preproc\_nest\_level} & \texttt{\$if} depth at include start; paired on EOF unwind. \\
\texttt{preproc\_word\_chars} & Identifier characters: letters, digits, \texttt{\_}. \\
\texttt{preproc\_command} & Current directive name, for diagnostics. \\
\texttt{preproc\_err\_count} & Errors reported. \\
\texttt{preproc\_dollar\_or\_pound} & Whether the directive started with \texttt{\#} or \texttt{\$}. \\
\texttt{preproc\_fmacro\_parse\_error} & Blocks object-like fallback after a failed \\
\texttt{preproc\_assert\_uid} & Suffix for unique temporaries in \texttt{assert} expansions. \\
\end{longtable}

\subsection{3.3 Driver}
\label{utr24-sec-driver}

\texttt{preprocessor()} is a generator. The translator drains it into \texttt{yyin}:

\begin{Verbatim}[fontsize=\small,frame=single,commandchars=\\\{\}]
every yyin ||:= preprocessor(fname, uni_predefs) do
   yyin ||:= "\\n"
\end{Verbatim}


in \texttt{uni/unicon/unicon.icn}. \texttt{uni\_predefs} is \texttt{predefs()} plus \texttt{-D} options. After the generator finishes, a non-zero \texttt{preproc\_err\_count} or a non-empty \texttt{parsingErrors} aborts before \texttt{yyparse()}.


Each line from \texttt{preproc\_read()} is scanned under \texttt{line~?~\{~...~\}}. After leading \texttt{preproc\_space()}:

\begin{itemize}
\item 
\texttt{\#line} (optional spaces, GCC form) goes to \texttt{preproc\_scan\_directive()}.
\item 
\texttt{\$} followed by a character in \texttt{nonpunctuation} (letters, digits, space, tab, form feed, carriage return) goes to \texttt{preproc\_scan\_directive()}.
\item 
Otherwise \texttt{\&pos~:=~1} and the line goes to \texttt{preproc\_scan\_text()}.
\end{itemize}

Unicon directives are \texttt{\$}-prefixed. A \texttt{\#} that is not \texttt{\#line} is not a directive. With macros loaded, \texttt{preproc\_scan\_text()} treats \texttt{\#} outside strings as a comment to end of line. Do not rely on C-style \texttt{\#define} at column 0.


There is no second pass over the same line.

\subsection{3.4 Conditional compilation}
\label{utr24-sec-ifstate}

\texttt{/preproc\_if\_state} is Unicon's null test. After \texttt{\$ifdef}, the true branch sets \texttt{preproc\_if\_state} to null (emit); the false branch sets \texttt{"false"} or \texttt{"off"} (skip). So \texttt{/preproc\_if\_state} means the current region is active (or there is no \texttt{\$ifdef}). \texttt{\textbackslash{}preproc\_if\_state} is the nonnull test.


\texttt{\$define}, \texttt{\$undef}, \texttt{\$include}, and text expansion run only when \texttt{/preproc\_if\_state} succeeds. \texttt{preproc\_scan\_text()} returns without emitting in skipped regions -- no triple-string rewrite, no macro expansion. \texttt{\$ifdef} / \texttt{\$ifndef} / \texttt{\$else} / \texttt{\$endif} still run in skipped regions so nested conditionals stay balanced. A bare \texttt{if} token inside skipped text is mapped to \texttt{\$if}, which pushes state and forces \texttt{"off"}.

\subsection{3.5 Text scan order}
\label{utr24-sec-text-order}

When the region is active, \texttt{preproc\_scan\_text()}:

\begin{enumerate}
\item 
Rewrites \texttt{"""} ... \texttt{"""} via \texttt{preproc\_rewrite\_triple\_strings()} before identifier scans, so the body is never treated as directives or macros.
\item 
If either symbol table is non-empty, scans for \texttt{\#}, strings, and identifiers. Otherwise it emits the line after \texttt{preproc\_sync\_lines()}.
\item 
Looks up function-like macros before object-like. A name in \texttt{preproc\_fun\_table} is a call only if \texttt{preproc\_expand\_fmacro\_call()} sees \texttt{(} after optional space; otherwise \texttt{\&pos} is restored and the name may still match an object-like macro.
\item 
Rescans replacements through \texttt{preproc\_scan\_text(done\_set)}. \texttt{done\_set} is \texttt{\&null}, a string, or a set of names on the current expansion chain; it blocks infinite recursion.
\end{enumerate}

String literals \texttt{"..."} / \texttt{'...'} are skipped with escape and trailing-\texttt{\_} continuation (\texttt{preproc\_read()}, adjusting \texttt{preproc\_print\_line}).

\section{4. Directives}
\label{utr24-sec-directives}

\texttt{preproc\_scan\_directive()} reads the command with \texttt{preproc\_word()}. Unknown names in an active region call \texttt{preproc\_error("unknown~preprocessor~directive")}.

\begin{longtable}{ll}
\hline
Command & Behavior \\
\hline
\endfirsthead
\hline
Command & Behavior \\
\hline
\endhead
\hline
\endfoot
\hline
\endlastfoot
\texttt{define} & Function-like \texttt{name(params...)~body} via \\
\texttt{undef} & Removes the name from both tables and from the builtin \\
\texttt{ifdef} / \texttt{ifndef} & Pushes \texttt{preproc\_if\_state}; sets state from \\
\texttt{\$if} & Placeholder for \texttt{if} tokens inside excluded code: \\
\texttt{else} / \texttt{endif} & Pop, toggle, or close the branch; error on mismatch. \\
\texttt{include} & Quoted or bare filename via \texttt{preproc\_qword()}; circular \\
\texttt{line} & Sets \texttt{preproc\_filename} / \texttt{preproc\_line} from a \\
\texttt{error} & User error; message to end of line or \texttt{\#}. \\
\texttt{ITRACE} & Sets \texttt{\&trace} from an integer. \\
\texttt{C} & Accumulates lines until \texttt{\$Cend}; passes the text to \\
\end{longtable}

\section{5. Triple-quoted multiline strings}
\label{utr24-sec-triple}

Delimiters \texttt{"""} ... \texttt{"""} mark a string whose body may span physical lines. The preprocessor collects raw text -- including lines that look like comments or \texttt{\$} directives -- and rewrites the region to an ordinary quoted literal. The lexer grammar in \texttt{uni/parser} is unchanged.

\begin{Verbatim}[fontsize=\small,frame=single,commandchars=\\\{\}]
procedure main()
   local x
   x := """line one
line two"""
   write(x)
end
\end{Verbatim}


Output:

\begin{Verbatim}[fontsize=\small,frame=single,commandchars=\\\{\}]
line one
line two
\end{Verbatim}


The opener may sit on the same line as code or alone on a following line. Indentation and newlines in the body are kept. Multiline bodies do not shift \texttt{\&line} for surrounding statements (\texttt{triple\_strings.icn}).

\begin{Verbatim}[fontsize=\small,frame=single,commandchars=\\\{\}]
procedure main()
   local x
   x := """
$ifdef __debug__
inside triple
$endif
"""
   write("[", x, "]")
   x := """escaped triple: \\"\\"\\" inside"""
   write(x)
   x := "contains \\"\\"\\" literally"
   write(x)
   x := """ends-with-backslash\\"""
   write("[", x, "]")
end
\end{Verbatim}


Output:

\begin{Verbatim}[fontsize=\small,frame=single,commandchars=\\\{\}]
[
$ifdef __debug__
inside triple
$endif
]
escaped triple: """ inside
contains """ literally
[ends-with-backslash\\]
\end{Verbatim}


Inside \texttt{"..."} or \texttt{'...'}, a \texttt{"""} sequence is copied through and does not start this mode. To put three double quotes in a triple-quoted body, write \texttt{\textbackslash{}"\textbackslash{}"\textbackslash{}"} (backslash before each quote). The four-character form \texttt{\textbackslash{}"""} is not an escape, so a body may end with a single \texttt{\textbackslash{}}.


\texttt{\$ifdef} / \texttt{\$endif} and \texttt{\#} comments inside the body are string content. Continuation lines pulled for an open triple are never scanned as directives.


\texttt{preproc\_rewrite\_triple\_strings()} runs at the start of \texttt{preproc\_scan\_text()}, only in an active region, and only if the current line contains \texttt{"""}. Outside quotes it still honors \texttt{\#} as a comment to end of that physical line. After an opener, it accumulates a body until the closer; if the line ends first, it appends a newline, \texttt{preproc\_read()}s the next line, and continues. An unclosed delimiter at EOF is \texttt{unterminated~triple-quoted~string}.


The body is turned into one ordinary string with \texttt{preproc\_quote\_string()}: named Unicon escapes for specials (including DEL as \texttt{\textbackslash{}d}), backslash and quote escaped, other C0 bytes as three-digit octal via \texttt{preproc\_octal\_escape()}. \texttt{triple\_strings.icn} checks that a DEL byte in a triple string survives compilation (\texttt{ord} 127).


Tests: \texttt{tests/unicon/triple\_strings.icn}, \texttt{tests/unicon/stand/triple\_strings.std}.

\section{6. Function-like macros}
\label{utr24-sec-fmacros}

Function-like macros extend object-like \texttt{\$define} with a parameter list and a replacement body. Nested calls, zero arguments, and commas inside quoted arguments are supported (\texttt{macros.icn}):

\begin{Verbatim}[fontsize=\small,frame=single,commandchars=\\\{\}]
$define SUM(a,b) ((a) + (b))
$define TWICE(x) SUM(x,x)
$define ZERO() 0
procedure main()
   write(SUM(2, 3))
   write(TWICE(5))
   write(ZERO())
end
\end{Verbatim}


Output:

\begin{Verbatim}[fontsize=\small,frame=single,commandchars=\\\{\}]
5
10
0
\end{Verbatim}


Quoted commas and parentheses do not split arguments. A call may continue on the next line without \texttt{\textbackslash{}}:

\begin{Verbatim}[fontsize=\small,frame=single,commandchars=\\\{\}]
$define ID(x) x
$define WRAP3(a,b,c) ((a) || "|" || (b) || "|" || (c))
$define SUM(a,b) ((a) + (b))
procedure main()
   write(ID("a,b"))
   write(WRAP3("a,b", "c(d)", "e"))
   write(SUM(1,
      2))
end
\end{Verbatim}


Output:

\begin{Verbatim}[fontsize=\small,frame=single,commandchars=\\\{\}]
a,b
a,b|c(d)|e
3
\end{Verbatim}


On \texttt{\$define}, \texttt{(} must follow the name immediately. \texttt{\$define~SUM~(a,b)~...} is object-like: the value starts at the space. At a call site, space before \texttt{(} is allowed (\texttt{preproc\_opt\_space()}), so a body may write \texttt{SUM~(x,y)}.


Zero parameters are legal (\texttt{ZERO()}). Whitespace-only between \texttt{(} and \texttt{)} is zero arguments (\texttt{ZERO(~)}), not one empty argument.


The replacement may continue on following lines if each continued line ends with \texttt{\textbackslash{}} as the very last character (no trim). \texttt{preproc\_scan\_define\_value()} and \texttt{preproc\_define\_value\_next\_line()} splice those lines: the \texttt{\textbackslash{}} and newline are dropped; other spaces are kept. A \texttt{\textbackslash{}} with trailing spaces is not a continuation. Leading space is skipped only before the first line of the body. The same \texttt{\textbackslash{}} rule applies to object-like \texttt{\$define} values (\texttt{macros\_define\_continuation.icn}). EOF after a splice is \texttt{unterminated~\$define~(line~continuation)}.

\begin{Verbatim}[fontsize=\small,frame=single,commandchars=\\\{\}]
$define GREETING "hello " || \\
"there"
$define SUMCONT(a,b) ((a) + \\
(b))
procedure main()
   write(GREETING)
   write(SUMCONT(2, 3))
end
\end{Verbatim}


Output:

\begin{Verbatim}[fontsize=\small,frame=single,commandchars=\\\{\}]
hello there
5
\end{Verbatim}


A call \texttt{NAME(arg1,~...)} may span physical lines without \texttt{\textbackslash{}}. \texttt{preproc\_scan\_macro\_args()} pulls more lines via \texttt{preproc\_macro\_arg\_next\_line()} until the closing \texttt{)}. A backslash inside a call is ordinary argument text (\texttt{macros\_call\_backslash\_literal.icn}), not a \texttt{\$define} splice.


Arguments split on top-level commas, with parenthesis depth, and with \texttt{"..."} / \texttt{'...'} opaque (commas and parens inside quotes do not split). \texttt{preproc\_subst\_fmacro()} replaces whole identifiers in the body, not substrings, and does not substitute inside strings. The result is rescanned, so \texttt{TWICE} and chains such as \texttt{A} -> \texttt{E} expand in one pipeline. \texttt{done\_set} blocks expanding the same name again on that chain.


Empty arguments (\texttt{SUM(1,)}) are \texttt{empty~macro~argument}. Duplicate parameter names are \texttt{duplicate~parameter~name~in~\$define}. Arity must match (\texttt{wrong~number~of~args~in~macro~call~to~NAME}). Unclosed \texttt{(} is \texttt{unterminated~macro~invocation}. If \texttt{(} follows the name on \texttt{\$define} but the parameter list is malformed, \texttt{preproc\_fmacro\_parse\_error} is set so the line is not treated as a bad object-like \texttt{\$define}.


Call continuation anchors \texttt{preproc\_print\_line} to the first physical line of the call (\texttt{preproc\_ml\_anchor}). \texttt{\&line} inside the expansion is that line; the next source line after the call skips the continuation rows. \texttt{\$define} \texttt{\textbackslash{}} splicing does not adjust \texttt{preproc\_print\_line}. \texttt{\&line} in a multiline body is the invocation line, not each physical line of the definition (\texttt{macros\_define\_multiline\_line.icn}).

\begin{Verbatim}[fontsize=\small,frame=single,commandchars=\\\{\}]
$define ASLINE(x) x

procedure main()
   write(&line)
   write(ASLINE(&line
   ))
   write(&line)
end
\end{Verbatim}


Output:

\begin{Verbatim}[fontsize=\small,frame=single,commandchars=\\\{\}]
4
5
7
\end{Verbatim}


(The call starts on line 5; line 6 is the continuation; the next statement is line 7.)


A function-like name is defined for \texttt{\$ifdef}. \texttt{\$undef} removes it. A name is either function-like or object-like, not both; \texttt{\$define} updates one table and clears the other. Redefinition of a function-like name is an error unless \texttt{preproc\_same\_fmacro()} finds the same parameters and body.

\begin{Verbatim}[fontsize=\small,frame=single,commandchars=\\\{\}]
$define SUM(a,b) ((a) + (b))
$ifdef SUM
$define HAS_SUM 1
$else
$define HAS_SUM 0
$endif
$undef SUM
$ifdef SUM
$define STILL 1
$else
$define STILL 0
$endif
procedure main()
   write(HAS_SUM)
   write(STILL)
end
\end{Verbatim}


Output:

\begin{Verbatim}[fontsize=\small,frame=single,commandchars=\\\{\}]
1
0
\end{Verbatim}


\texttt{\$undef~assert} then \texttt{\$define~assert(...)} replaces the builtin; \texttt{macros.icn} does this for both \texttt{assert} and \texttt{assert\_not}.


Tests: \texttt{tests/unicon/macros.icn}; \texttt{macros\_multiline\_line.icn}, \texttt{macros\_define\_multiline\_line.icn}, \texttt{macros\_define\_continuation.icn}, \texttt{macros\_call\_backslash\_literal.icn}; negatives under \texttt{tests/unicon/data/macros\_bad\_*.icn} and \texttt{macros\_bad\_define\_eof*.icn}.

\section{7. Built-in \texttt{assert} and \texttt{assert\_not}}
\label{utr24-sec-assert}

\texttt{preproc\_install\_builtin\_macros()} registers \texttt{assert} and \texttt{assert\_not} in \texttt{preproc\_fun\_table} as placeholder \texttt{preproc\_fmacro} records and marks them in \texttt{preproc\_builtin\_fun\_table}. Expansion is not \texttt{preproc\_subst\_fmacro()}: \texttt{preproc\_expand\_fmacro\_call()} builds a Unicon fragment from the argument text. Arity is one or two: \texttt{assert(expr)} or \texttt{assert(expr,~label)}. Any other count is \texttt{wrong~number~of~args~in~macro~call}.


Semantics follow Unicon success and failure, not Boolean truth: \texttt{assert(expr)} expects \texttt{expr} to succeed; \texttt{assert\_not(expr)} expects \texttt{expr} to fail. \texttt{assert} is alternation (\texttt{expr~|~\{~...~\}}); \texttt{assert\_not} is \texttt{if~(expr)~then~\{~...~\}}.

\begin{longtable}{lll}
\hline
Mode & Enable & On failure \\
\hline
\endfirsthead
\hline
Mode & Enable & On failure \\
\hline
\endhead
\hline
\endfoot
\hline
\endlastfoot
Debugging & \texttt{\$define~\_\_debug\_\_} or \texttt{-D\_\_debug\_\_} & \texttt{runerr(219,~payload)}. The process stops. Traceback \\
Testing & \texttt{\$define~\_\_test\_\_} or \texttt{-D\_\_test\_\_} & \texttt{write} of \texttt{[file:line]~assertion~failed~(expr)} \\
Release & neither & Expands to nothing (\texttt{assert\_strip.icn}). \\
\end{longtable}


If both \texttt{\_\_test\_\_} and \texttt{\_\_debug\_\_} are defined, the testing path (\texttt{fail}) wins. \texttt{preproc\_new()} and a later \texttt{\$define~\_\_test\_\_} both default \texttt{\_\_debug\_\_} when it is absent.


These macros are standalone statements. Disabled expansion is empty, so expression positions such as \texttt{if~assert(...)} are not supported.


The debug payload is the quoted expression text, then \texttt{char(30)} and the label if the label is non-empty (\texttt{AssertRunerrSep} in \texttt{errmsg.r}). \texttt{assert\_not} prefixes \texttt{char(29)} so the traceback can print \texttt{assert\_not(}. Error 219 is \texttt{assertion~failed} (\texttt{data.r}). \texttt{errmsg.r} skips the ordinary \texttt{Run-time~error~219} banner and the offending- value line; it prints \texttt{[file:line]~assertion~failed:~...}.


If the label expression fails, the assertion diagnostic still runs: the expansion does \texttt{lbl~:=~((label)~|~"")} (\texttt{assert\_failing\_label.icn}).


Each expansion increments \texttt{preproc\_assert\_uid} so generated temporaries (\texttt{\_\_preproc\_assert\_N\_lbl}) do not collide -- icont has no block-local variables in expression-level \texttt{\{...\}} compounds. The fragment is rescanned so nested macros in the inserted code expand.


\texttt{\$undef~assert} then \texttt{\$define~assert(...)} replaces the builtin like any user function-like macro.

\begin{Verbatim}[fontsize=\small,frame=single,commandchars=\\\{\}]
$define __debug__
procedure main()
   local x
   x := 1
   assert(x = 1)
   assert_not(x = 2)
   assert(2 === 1 + 1, "1 + 1 = 2")
   write("ok")
end
\end{Verbatim}


Output:

\begin{Verbatim}[fontsize=\small,frame=single,commandchars=\\\{\}]
ok
\end{Verbatim}


A failing debug assert stops the process. The traceback names \texttt{assert} / \texttt{assert\_not} rather than \texttt{runerr}:

\begin{Verbatim}[fontsize=\small,frame=single,commandchars=\\\{\}]
$define __debug__
procedure main()
   assert_not(1 = 1)
   write("should not reach here")
end
\end{Verbatim}


Output:

\begin{Verbatim}[fontsize=\small,frame=single,commandchars=\\\{\}]
[assert_ex.icn:3] assertion failed: 1 = 1
Traceback:
   main()
   assert_not(1 = 1) from line 3 in assert_ex.icn
\end{Verbatim}


Under \texttt{\_\_test\_\_}, failure writes a diagnostic and \texttt{fail}s, so later statements still run:

\begin{Verbatim}[fontsize=\small,frame=single,commandchars=\\\{\}]
$define __test__
procedure main()
   assert(1 = 2, "keep going")
   write("after fail")
   assert(1 = 1)
   write("done")
end
\end{Verbatim}


Output:

\begin{Verbatim}[fontsize=\small,frame=single,commandchars=\\\{\}]
[assert_ex.icn:3] assertion failed (1 = 2) keep going
after fail
done
\end{Verbatim}


With neither symbol, the calls vanish. The asserts below would fail at run time if they were still present, so reaching \texttt{write()} is the proof they were stripped:

\begin{Verbatim}[fontsize=\small,frame=single,commandchars=\\\{\}]
procedure main()
   local x
   x := 1
   assert(x = 2)
   assert_not(x = 1)
   write("strip-ok")
end
\end{Verbatim}


Output:

\begin{Verbatim}[fontsize=\small,frame=single,commandchars=\\\{\}]
strip-ok
\end{Verbatim}


Tests: \texttt{assert\_debugging.icn} (passing cases, then a failing assert five frames down), \texttt{assert\_testing.icn}, \texttt{assert\_strip.icn}, \texttt{assert\_failing\_label.icn}, \texttt{assert\_not\_failing.icn}, matching \texttt{stand/*.std}. Builtin override is in \texttt{macros.icn}.

\begin{longtable}{ll}
\hline
Concern & Primary routines (\texttt{uni/unicon/preproce.icn}) \\
\hline
\endfirsthead
\hline
Concern & Primary routines (\texttt{uni/unicon/preproce.icn}) \\
\hline
\endhead
\hline
\endfoot
\hline
\endlastfoot
Triple-quoted strings & \texttt{preproc\_rewrite\_triple\_strings()}, \\
Function-like macros & \texttt{record~preproc\_fmacro}, \\
Assertions & \texttt{preproc\_install\_builtin\_macros()}, \\
\end{longtable}

\section{8. Line numbers, errors, and includes}
\label{utr24-sec-support}

\texttt{preproc\_sync\_lines()} emits \texttt{\#line~N~"file"} when the filename changes or the gap is large; for small gaps it may emit blank lines. \texttt{preproc\_scan\_text()} can \texttt{suspend} more than once per logical line. Features that insert, delete, or merge lines must keep \texttt{preproc\_line} and \texttt{preproc\_print\_line} consistent with that.


\texttt{preproc\_error(msg)} builds \texttt{File~...;~Line~N~\#~\$directive:~...}, increments \texttt{preproc\_err\_count}, and pushes a \texttt{ParseError} onto \texttt{parsingErrors}.


\texttt{preproc\_read()} increments \texttt{preproc\_line} and returns the next line from \texttt{read()} or from a list "file." On EOF it checks \texttt{preproc\_if\_stack} depth against \texttt{preproc\_nest\_level}, closes the file, pops \texttt{preproc\_file\_stack}, and removes the path from \texttt{preproc\_include\_set}. String continuation, triple-quoted bodies, and macro arguments that need another physical line all call \texttt{preproc\_read()} again.


\texttt{preproc\_quote\_string(s)} escapes arbitrary text for a Unicon string literal. \texttt{preproc\_istring()} / \texttt{preproc\_istring\_radix()} decode escapes in quoted filenames.

\section{9. Extending the preprocessor}
\label{utr24-sec-extend}

All of the following is in \texttt{uni/unicon/preproce.icn} unless noted. Coordinate any change with \texttt{uni/parser/preproce.icn} if that copy is still shipped; it does not currently match.

\subsection{9.1 Kind and touchpoints}
\label{utr24-sec-extend-kinds}
\begin{longtable}{ll}
\hline
Kind & Touchpoints \\
\hline
\endfirsthead
\hline
Kind & Touchpoints \\
\hline
\endhead
\hline
\endfoot
\hline
\endlastfoot
New or changed \texttt{\$} directive & \texttt{preproc\_scan\_directive()}: add a \texttt{case} label; set \\
Object-like macro behavior & \texttt{preproc\_sym\_table}, substitution in \\
Function-like macro behavior & \texttt{preproc\_fun\_table}, \\
Built-in special function macro & \texttt{preproc\_install\_builtin\_macros()}, \\
Source syntax rewritten to plain Unicon & A phase in \texttt{preproc\_scan\_text()}. Must respect \\
New predefined symbol & \texttt{predefs()}. Symbols that depend on compiler options \\
\end{longtable}

\subsection{9.2 Guards and line accounting}
\label{utr24-sec-extend-guards}

Directives that define, undefine, or include use \texttt{if~/preproc\_if\_state~then} so they do not run when code is stripped. Directives that only keep parser state in skipped regions follow \texttt{ifdef} / \texttt{\$if} / \texttt{else} / \texttt{endif}, not that guard.


If the feature reads extra lines or changes output length, check \texttt{preproc\_line}, \texttt{preproc\_print\_line}, and \texttt{preproc\_sync\_lines()}. \texttt{unicon~-E} shows the emitted \texttt{\#line} stream. Multiple \texttt{suspend}s per input line must leave \texttt{preproc\_print\_line} consistent at the end of \texttt{preproc\_scan\_text()}.

\subsection{9.3 Errors and tests}
\label{utr24-sec-extend-tests}

Use \texttt{preproc\_error(msg)} after \texttt{preproc\_command} is set so the message includes \texttt{\$directive:}. The compiler exit path is \texttt{preproc\_err\_count} and \texttt{parsingErrors} in \texttt{unicon.icn}.


Add a driver under \texttt{tests/unicon/} and expected output under \texttt{stand/}. Error cases belong in \texttt{tests/unicon/data/}. Update \hyperref[utr24-sec-triple]{sections 5--7} when user-visible semantics change.

\subsection{9.4 New \texttt{\$} directive}
\label{utr24-sec-extend-directive}

After \texttt{preproc\_command~:=~preproc\_word()}, the \texttt{case~preproc\_command~of} dispatch matches the word after \texttt{\$}. A new label:

\begin{Verbatim}[fontsize=\small,frame=single,commandchars=\\\{\}]
"mydir": \{
   if /preproc_if_state then \{
      # parse remainder of line; on failure:
      preproc_error("optional detail")
   \}
\}
\end{Verbatim}


Unknown names still hit \texttt{default} and \texttt{preproc\_error("unknown~preprocessor~directive")} when active.

\section{10. Related documents}
\label{utr24-sec-related}
\begin{longtable}{ll}
\hline
Document & Content \\
\hline
\endfirsthead
\hline
Document & Content \\
\hline
\endhead
\hline
\endfoot
\hline
\endlastfoot
UTR \#15 [1] & How to write and submit a Unicon Technical Report. \\
UTR \#8 [2] & Language reference preprocessor commands \\
Implementation book [4] & \texttt{doc/ib/p3-unicon.tex}, "The Unicon Preprocessor": \\
\end{longtable}


Other \texttt{doc/ib/} preprocessor mentions (Idol, Ratfor, rtt) are different tools, not \texttt{preproce.icn}.

\section{References}
\label{utr24-references}
\begin{thebibliography}{99}
\bibitem[1]{cite-jeffery-utr15}
Clinton Jeffery. "How to Write a Unicon Technical Report". doc/utr/utr15.tex. 2013.

\bibitem[2]{cite-jeffery-utr8}
Clinton Jeffery. "Unicon Language Reference". doc/unicon/utr8.html.

\bibitem[3]{cite-alexander-jcon-preproc}
Bob Alexander. "Preprocessor for Jcon". Basis for Unicon's preproce.icn.

\bibitem[4]{cite-jeffery-ib-preproc}
Clinton Jeffery. "The Unicon Compiler". Course notes, doc/ib/p3-unicon.tex, section "The Unicon Preprocessor".

\bibitem[5]{cite-libssh}
The libssh Project. "libssh -- The SSH Library". https://www.libssh.org/. 2026.

\bibitem[6]{cite-jeffery-utr13}
Clinton Jeffery. "The Unicon Messaging Facilities". doc/utr/utr13.odt.

\bibitem[7]{cite-openssl}
The OpenSSL Project. "OpenSSL -- Cryptography and SSL/TLS Toolkit". https://www.openssl.org/. 2026.

\bibitem[8]{cite-algharaibeh-utr26}
Jafar Al-Gharaibeh. "Native SSH and SFTP Support in Unicon". doc/utr/utr26.rst. 2026.

\bibitem[9]{cite-rfc2104}
H. Krawczyk and M. Bellare and R. Canetti. "HMAC: Keyed-Hashing for Message Authentication". RFC 2104. 1997.

\end{thebibliography}
\end{document}
